% Part: computability
% Chapter: recursive-functions
% Section: other-recursions

\documentclass[../../../include/open-logic-section]{subfiles}

\begin{document}

\olfileid{cmp}{rec}{ore}
\olsection{عوديات أخرى}

يمكننا، باستعمال الاقتران والترتيب في متتاليات، تسويغ صور أغرب (وأنفع)
من العودية البدائية. فكثيرًا ما يفيد، مثلًا، تعريف دالتين في آن واحد،
كما في التعريف الآتي:
\begin{align*}
h_0(\vec x, 0) & = f_0(\vec x) \\
h_1(\vec x, 0) & = f_1(\vec x) \\
h_0(\vec x, y+1) & = g_0(\vec x, y, h_0(\vec x, y), h_1(\vec x, y)) \\
h_1(\vec x, y+1) & = g_1(\vec x, y, h_0(\vec x, y), h_1(\vec x, y))
\end{align*}
وهذه حالة من \emph{العودية المتزامنة}. ومن الطرق المفيدة الأخرى لتعريف
الدوال إعطاء قيمة $h(\vec x,y+1)$ بدلالة \emph{جميع} القيم
$h(\vec x,0)$، …،~$h(\vec x,y)$، كما في التعريف الآتي:
\begin{align*}
h(\vec x, 0) & = f(\vec x) \\
h(\vec x, y+1) & = g(\vec x, y, \tuple{h(\vec x, 0), \dots, h(\vec x, y)}).
\end{align*}
ويصور المخطط الآتي هذه الفكرة بإيجاز أكبر:
\[
h(\vec x, y) = g(\vec x, y, \tuple{h(\vec x, 0), \dots, h(\vec x, y-1)})
\]
على أن نفهم أن الحجة الأخيرة لـ$g$ هي المتتالية الخالية حين يكون
$y=0$. والفكرة في كلتا الصورتين هي أن الدالة $h$ تستطيع، عند حساب
«خطوة اللاحق»، استعمال متتالية القيم المحسوبة حتى ذلك الموضع بأكملها.
وتعرف هذه باسم \emph{عودية مسار القيم}. ويمكن استعمالها، في مثال
معين، لتسويغ تعريف من النوع الآتي:
\begin{align*}
h(\vec x, y) & = \begin{cases}
  g(\vec x, y, h(\vec x, k(\vec x, y))) & \text{إذا كان $k(\vec x, y) < y$} \\
  f(\vec x) & \text{في غير ذلك}
\end{cases}
\end{align*}
أي إن قيمة $h$ عند $y$ يمكن حسابها بدلالة قيمة $h$ عند
\emph{أي} قيمة سابقة، تحددها~$k$.

\begin{prob}
  عرف دالة الباقي $r(x,y)$ بعودية مسار القيم. (إذا كان $x$ و$y$
  عددين طبيعيين وكان $y>0$، فإن $r(x,y)$ هو العدد الأصغر من~$y$ الذي
  يحقق $x=z\times y+r(x,y)$ لبعض~$z$. وللتعيين، لنقل إنه إذا كان
  $y=0$ فإن $r(x,0)=0$.)
\end{prob}

ينبغي أن تفكر في كيفية الحصول على هذه الدوال بالعودية البدائية
المعتادة. وثمة صورة أخيرة من العودية البدائية أشد مرونة، إذ يسمح فيها
بتغيير \emph{المعلمات} (القيم الجانبية) على امتداد الطريق:
\begin{align*}
h(\vec x, 0) & = f(\vec x) \\
h(\vec x, y+1) & = g(\vec x, y, h(k(\vec x), y))
\end{align*}
ويمكن محاكاة هذه أيضًا بالعودية البدائية المعتادة. (والقيام بذلك دقيق؛
وكتلميح، حاول بسط الحساب يدويًا.)

\end{document}
